\documentclass[11pt,a4paper]{article}

% ---------- Packages ----------
\usepackage[margin=2.2cm]{geometry}
\usepackage{amsmath,amssymb,amsthm}
\usepackage{mathtools}
\usepackage{enumitem}
\usepackage{hyperref}
\usepackage{ctex} % for Chinese; if you don't need Chinese, remove this line

\setlist[itemize]{leftmargin=1.6em}
\setlist[enumerate]{leftmargin=1.8em}

% ---------- Macros ----------
\newcommand{\R}{\mathbb{R}}
\newcommand{\Om}{\Omega}
\newcommand{\dd}{\,\mathrm{d}}
\newcommand{\dx}{\mathrm{d}x}
\newcommand{\dy}{\mathrm{d}y}
\newcommand{\df}{\mathrm{d}f}
\newcommand{\pull}{^{*}} % pullback star

\title{微分形式（Differential Forms）常用运算法则速查表}
\author{}
\date{}

\begin{document}
\maketitle

\section*{0.\ 基本对象与记号}
设 $U\subset \R^n$ 为开集，$T_pU\simeq \R^n$。
\begin{itemize}
  \item \textbf{0-形式}：可微函数 $g:U\to\R$，记作 $g\in \Om^0(U)$。
  \item \textbf{$k$-形式}：对每个 $p\in U$ 赋予一个交替 $k$-线性映射
  \[
  \omega(p):(T_pU)^k\to \R,\qquad \omega\in \Om^k(U).
  \]
  \item \textbf{坐标展开}：
  \[
  \omega=\sum_{I} a_I\, \dx_I,\qquad
  \dx_I=\dx_{i_1}\wedge\cdots\wedge \dx_{i_k},\ i_1<\cdots<i_k.
  \]
  \item \textbf{维数限制}：若 $k>n$，则 $\Om^k(U)=\{0\}$。
\end{itemize}

\section*{1.\ 加法与函数乘法（0-形式乘 $k$-形式）}
\begin{itemize}
  \item \textbf{加法（同阶）}：$(\omega+\eta)(p)=\omega(p)+\eta(p)$。
  \item \textbf{函数乘形式（按点标量乘法）}：
  \[
  (g\omega)(p)(v_1,\dots,v_k)=g(p)\,\omega(p)(v_1,\dots,v_k).
  \]
  \item 也可视作 $g\wedge \omega$（因为 $g$ 是 0-形式）。
\end{itemize}

\section*{2.\ 楔积（外积）$\wedge$}
若 $\omega\in\Om^k(U)$，$\eta\in\Om^s(U)$，则 $\omega\wedge\eta\in\Om^{k+s}(U)$。

\subsection*{2.1\ 双线性与分配律}
\[
(\omega_1+\omega_2)\wedge\eta=\omega_1\wedge\eta+\omega_2\wedge\eta,\qquad
\omega\wedge(\eta_1+\eta_2)=\omega\wedge\eta_1+\omega\wedge\eta_2.
\]
若 $g,h\in\Om^0(U)$，
\[
(g\omega)\wedge(h\eta)=(gh)(\omega\wedge\eta).
\]

\subsection*{2.2\ 结合律}
\[
(\omega\wedge\eta)\wedge\theta=\omega\wedge(\eta\wedge\theta).
\]

\subsection*{2.3\ 分次交换律（graded commutativity）}
若 $\deg\omega=k$，$\deg\eta=s$，则
\[
\omega\wedge\eta = (-1)^{ks}\,\eta\wedge\omega.
\]
推论：若 $k$ 为奇数，则 $\omega\wedge\omega=0$。

\subsection*{2.4\ 基本元取值（常用）}
若 $i_1<\cdots<i_k$ 且 $j_1<\cdots<j_k$，则
\[
(\dx_{i_1}\wedge\cdots\wedge \dx_{i_k})(e_{j_1},\dots,e_{j_k})
=
\begin{cases}
1,& (i_1,\dots,i_k)=(j_1,\dots,j_k),\\
0,& \text{否则}.
\end{cases}
\]

\section*{3.\ 外微分 $d$}
外微分是线性算子 $d:\Om^k(U)\to \Om^{k+1}(U)$。

\subsection*{3.1\ 对 0-形式（函数）}
若 $g\in\Om^0(U)$，
\[
dg=\sum_{i=1}^n \frac{\partial g}{\partial x_i}\,\dx_i.
\]

\subsection*{3.2\ 对一般 $k$-形式（书中定义）}
若 $\omega=\sum_I a_I\,\dx_I$，则
\[
d\omega=\sum_I da_I\wedge \dx_I.
\]

\subsection*{3.3\ 线性}
\[
d(\omega+\eta)=d\omega+d\eta.
\]

\subsection*{3.4\ Leibniz 公式（分次乘法法则）}
若 $\deg\omega=k$，则
\[
d(\omega\wedge\eta)=d\omega\wedge\eta+(-1)^k\,\omega\wedge d\eta.
\]
特别地，若 $g$ 为 0-形式，
\[
d(g\omega)=dg\wedge\omega+g\,d\omega.
\]

\subsection*{3.5\ $d^2=0$}
\[
d(d\omega)=0.
\]
特别地 $d(\dx_i)=0$，因此 $d(\dx_I)=0$。

\section*{4.\ 拉回（变量代换）$f\pull$}
设 $f:U\subset\R^n\to V\subset\R^m$ 可微，则
\[
f\pull:\Om^k(V)\to \Om^k(U).
\]
定义（在点 $p$）：
\[
(f\pull\omega)(p)(v_1,\dots,v_k)=\omega(f(p))\big(df_p(v_1),\dots,df_p(v_k)\big).
\]

\subsection*{4.1\ 对 0-形式（函数）}
\[
f\pull g=g\circ f.
\]

\subsection*{4.2\ 线性与函数乘法相容}
\[
f\pull(\omega+\eta)=f\pull\omega+f\pull\eta,\qquad
f\pull(g\omega)=f\pull(g)\,f\pull(\omega).
\]
（右边乘法是“函数 $\times$ 形式”的按点标量乘法。）

\subsection*{4.3\ 与楔积相容}
\[
f\pull(\omega\wedge\eta)=f\pull\omega\wedge f\pull\eta.
\]

\subsection*{4.4\ 复合律}
若 $g:W\to U$ 可微，则
\[
(f\circ g)\pull = g\pull\circ f\pull.
\]

\subsection*{4.5\ 与外微分可交换}
\[
d(f\pull\omega)=f\pull(d\omega).
\]

\subsection*{4.6\ 维数原因导致为 0}
若 $\dim U=m$ 且 $k>m$，则对任意 $\omega\in\Om^k(V)$，
\[
f\pull\omega=0.
\]

\section*{5.\ 坐标计算速记（代换规则）}
若 $f=(f_1,\dots,f_m)$，即 $y_i=f_i(x)$，则
\[
f\pull(\dy_i)=d(y_i\circ f)=df_i=\sum_{j=1}^n \frac{\partial f_i}{\partial x_j}\,\dx_j.
\]
因此对
\[
\omega=\sum_I a_I(y)\,\dy_I
\]
有
\[
f\pull\omega=\sum_I a_I(f(x))\, df_{i_1}\wedge\cdots\wedge df_{i_k}.
\]
（把 $y$ 换成 $f(x)$，把 $\dy_i$ 换成 $df_i$。）

\section*{6.\ 内积诱导的向量场与 1-形式对应}
在 $\R^n$ 标准内积 $\langle\cdot,\cdot\rangle$ 下：
\[
v\in T_p\R^n \longmapsto v^\flat\in T_p^*\R^n,\qquad
v^\flat(u)=\langle v,u\rangle.
\]
这给出向量场与 1-形式的逐点一一对应（常称 ``$\flat/\sharp$'' 音乐同构）。

\vspace{0.6em}
\hrule
\vspace{0.6em}
\noindent\textbf{备注：}本文未涉及 Hodge 星号 $\star$（那是依赖度量的 $k\leftrightarrow (n-k)$ 运算，与对偶/拉回的 $*$ 不同）。

\end{document}
